%% ==================================================
%% Appendix: supporting results for the physics examples of the
%% multi-block compression theorem: explicit gauges, adapted bases,
%% and a sign-weighted pair used as a test construction.
%% ==================================================

\chapter{Physics Examples of the Multi-Block Compression Theorem: Supporting Results}%
\label{ch:asymmetric_examples_supporting}

This appendix supports Chapter~\ref{ch:asymmetric_examples}. It collects the
explicit gauges and the integer arithmetic behind the compression theorems of
the examples there that are treated in explicit bond coordinates, the adapted
bases and scaled matrix units behind the normality certificates and behind the
compression of the Bell-bond fixed points, and a sign-weighted pair of blocks
constructed in this development to test the sharpness of the compression
theorem.

\section{Explicit gauges and exact arithmetic}%
\label{sec:asymex_explicit_gauge}

This section supplies the change of bond coordinates and the integer
certificates used by the compression theorems of the examples in
Chapter~\ref{ch:asymmetric_examples} that are treated in explicit bond
coordinates: the symmetry and fusion operators, the elementary states, and
the sign-weighted pair of Section~\ref{sec:asymex_parity}.
Theorem~\ref{thm:asymex_czx_compression} and
Theorem~\ref{thm:asymex_ising_compression}, for the CZX operator and the
Ising fusion operators of the review~\cite{Cirac2021Matrix}, invoke it
directly; the Bell-bond fixed points rest on
Section~\ref{sec:asymex_rfp_basis} instead.
The compression datum of Theorem~\ref{thm:asym_multiblock}
carries its change of bond coordinates as an invertible map onto a graded
coordinate space, one summand per slot. In each worked example that map is
an explicit invertible matrix together with an explicit labelling of the
bond coordinates by slots and positions inside slots.
Theorem~\ref{thm:asymex_explicit_gauge} reads the compression pairs of a
slot off that matrix, Theorem~\ref{thm:asymex_one_slot} treats the
recurring one-slot shape, and Theorem~\ref{thm:asymex_int_certificates}
turns integer identities into certificates of normality and of vanishing
sitewise intertwiners.

\begin{definition}[Explicit gauge and entrywise integer coercion]
    \label{def:asymex_explicit_gauge}
    \lean{MPSTensor.complexOfInt, MPSTensor.mulIntTensor, MPSTensor.gaugeOfMatrix}
    \leanok
    \uses{def:mps_tensor, def:mpdo_operator_product}
    Let $\tau$ be a bijection between the graded coordinate space of a slot
    family and the bond index set $\{0,\ldots,D_B-1\}$, and let $G$ and
    $G^{-1}$ be mutually inverse matrices in $\MN{D_B}$. The \emph{explicit
        gauge} attached to $(\tau,G)$ is the invertible map sending a bond
    vector $v$ to the graded vector $b\mapsto(Gv)(\tau b)$. For an integer
    matrix $X$, its \emph{entrywise coercion} is the complex matrix with the
    same entries; for integer matrix product operator tensors $M$ and $N$,
    their integer bond-space product is
    $(M\mathbin{\cdot}N)^{ik}=\sum_jM^{ij}\otimes N^{jk}$.
\end{definition}

\begin{theorem}[Conjugation, coercion and the compression pairs of a slot]
    \label{thm:asymex_explicit_gauge}
    % Principal declaration: MPSTensor.conjMatrix_gaugeOfMatrix.
    \lean{MPSTensor.conjMatrix_gaugeOfMatrix, MPSTensor.complexOfInt_mul,
        MPSTensor.mulTensor_smul_complexOfRing,
        MPSTensor.MultiBlockCompression.left_gaugeOfMatrix,
        MPSTensor.MultiBlockCompression.right_gaugeOfMatrix}
    \leanok
    \uses{def:asymex_explicit_gauge, thm:asym_multiblock}
    Conjugation by the explicit gauge attached to $(\tau,G)$ sends a bond
    matrix $A$ to the matrix $GAG^{-1}$ relabelled along $\tau$. Entrywise
    coercion is multiplicative, and for integer tensors $M, N$ and $c \in \C$,
    \begin{align}
        (cM) \mathbin{\cdot} (cN) = c^2 (M \mathbin{\cdot} N). \notag
    \end{align}
    If a compression datum uses the explicit gauge $(\tau,G)$, then the
    compression maps for slot $s$ are given in coordinates by
    $W_s = \pi_s G$ and $V_s = G^{-1} \iota_s$, where $\pi_s$ selects the rows
    and $\iota_s$ embeds the columns belonging to slot $s$.
\end{theorem}

\begin{proof}
    \leanok
    The matrix of the explicit gauge is $G$ relabelled along $\tau$ in its
    rows, and the matrix of its inverse is $G^{-1}$ relabelled along $\tau$
    in its columns, so conjugation is the stated matrix conjugation followed
    by the relabelling. Multiplicativity of the coercion is the fact that a
    sum of products of integers is the same computed in $\Z$ or in $\C$, and
    the statement about the bond-space product follows by expanding both
    sides entrywise. The two compression maps are the composites of the
    coordinate embedding and projection of the slot with the gauge and its
    inverse; evaluating the matrix products at a single graded coordinate
    leaves one row of $G$ and one column of $G^{-1}$.
\end{proof}

\begin{definition}[One-slot compression data]%
    \label{def:asymex_one_slot}
    \lean{MPSTensor.stackedInt, MPSTensor.oneSlot, MPSTensor.oneSlotMem,
        MPSTensor.MultiBlockCompression.ofConjMatrix,
        MPSTensor.MultiBlockCompression.ofConjInt,
        MPSTensor.MultiBlockCompression.ofEq,
        MPSTensor.MultiBlockCompression.ofScalarFlagFour}
    \leanok
    \uses{def:asymex_explicit_gauge, thm:asym_multiblock}
    For a stacked product of two integer matrix product operator tensors,
    read over the pair alphabet, the letter at $a=do+i$ is
    $(M\mathbin{\cdot}N)^{oi}$. A \emph{one-slot compression datum} for a
    tensor $B$ onto a single target $C$ is the data of
    Theorem~\ref{thm:asym_multiblock} for the one-element slot set: an
    explicit gauge $(\tau,G)$ conjugating $B$ into a matrix that is block
    upper triangular for a flag with $C$ at one diagonal position. Two
    shapes recur in the examples below: a conjugation $GB^iG^{-1}=C^i$ for
    an invertible $G$ gives such a datum with no zero slots ($z=0$); and,
    for a four-dimensional $B$ and a one-dimensional $C$, an integer gauge
    sending every letter to a matrix with $C$ at the third diagonal
    position of a four-step flag and zero at the other three gives such a
    datum with three zero slots ($z=3$).
\end{definition}

\begin{theorem}[The remainder of a conjugate datum vanishes; word traces
        of a one-slot datum]%
    \label{thm:asymex_one_slot}
    \lean{MPSTensor.MultiBlockCompression.remainder_ofConjMatrix,
        MPSTensor.MultiBlockCompression.remainder_ofEq,
        MPSTensor.MultiBlockCompression.remainder_ofConjInt,
        MPSTensor.MultiBlockCompression.trace_evalWord_oneSlot,
        MPSTensor.MultiBlockCompression.trace_evalWord_oneSlot_smul}
    \leanok
    \uses{def:asymex_one_slot}
    The remainder of a conjugate one-slot datum vanishes identically: the
    compression is an exact change of coordinates. For every one-slot
    datum with target tensor $C$ and nonempty word $w \ne \emptyset$,
    \begin{align}
        \tr(B^w) = \tr(C^w), \notag
    \end{align}
    or $\tr(B^w) = c^{|w|}\tr(C^w)$ when the target is carried with scalar
    weight $c$.
\end{theorem}

\begin{proof}
    \leanok
    \uses{prop:asym_compression_trace}
    Conjugation preserves the trace, which decomposes as the sum of traces of
    the diagonal blocks. For the upper block-triangular gauge, the diagonal
    blocks of a word product are the word products of the individual diagonal
    blocks; the zero-slot diagonal block is strictly upper triangular (or zero)
    and therefore has vanishing trace on any nonempty word $w \ne \emptyset$,
    leaving only the contribution of the target block $C$. The scalar factor
    carries across words by the length $|w|$.
\end{proof}

\begin{theorem}[Integer certificates for normality and for vanishing
        sitewise intertwiners]%
    \label{thm:asymex_int_certificates}
    \lean{MPSTensor.isNBlkInjective_two_of_int,
        MPSTensor.right_intertwiner_eq_zero_of_int,
        MPSTensor.left_intertwiner_eq_zero_of_int}
    \leanok
    \uses{def:normal}
    Let $A \in (\MN{D})^d$ have integer entries $A_i \in \mathrm{M}_D(\Z)$.
    If there exist index pairs $(a_k, b_k)$, integer coefficients
    $T_{xyk} \in \Z$, and a nonzero integer $m \ne 0$ such that for all matrix
    units $E_{xy}$,
    \begin{align}
        \sum_k T_{xyk} A_{a_k} A_{b_k} = m E_{xy}, \notag
    \end{align}
    then $\spn_\C\{A_i A_j : i,j \in \{0,\dots,d-1\}\} = \MN{D}$.
    Similarly, let $B \in (\mathrm{M}_D(\Z))^d$ and $c \in (\mathrm{M}_1(\Z))^d$.
    If a selection of rows from the matrices $B_i - c_i \Id_D$ forms an
    integer matrix $K \in \mathrm{M}_D(\Z)$ admitting an integer left inverse
    $K' K = m \Id_D$ with $m \ne 0$, then every $v \in \C^D$ satisfying
    $B_i v = c_i v$ for all $i$ satisfies $v = 0$.
    Likewise, if a selection of columns from $B_i - c_i \Id_D$ admits an
    integer left inverse up to $m \ne 0$, then every $u \in \C^D$ with
    $u B_i = c_i u$ satisfies $u = 0$.
\end{theorem}

\begin{proof}
    \leanok
    Dividing the linear combination by $m \ne 0$ exhibits each matrix unit
    $E_{xy}$ in the span of length-two words of $A$, showing they span
    $\MN{D}$. For the sitewise intertwiner, $B_i v = c_i v$ implies $K v = 0$,
    whence $m v = K' K v = 0$, so $v = 0$ since $m \ne 0$ in $\C$. The column
    argument for $u$ is identical.
\end{proof}

\section{Adapted bases and scaled matrix units}%
\label{sec:asymex_rfp_basis}

This section supplies two tools. Theorem~\ref{thm:asymex_normal_units}
certifies normality of a tensor whose letters contain a scaled copy of
every matrix unit; it is invoked in the CZX and Ising sections of
Chapter~\ref{ch:asymmetric_examples} and in the fixed-point section
there. Theorem~\ref{thm:asymex_stacked_pair} converts an integer matrix
identity, written in an adapted basis of the inner bond, into the weighted
direct-sum hypothesis of Proposition~\ref{prop:asym_projector_weighted_sum};
it is the form in which the vertical components of the renormalization
fixed points of Chapter~\ref{ch:asymmetric_examples} enter the compression
theorem.

The stacked bond spaces of those vertical components are tensor squares.
Reordering the four bond bits separates the two middle bits, which form
the inner bond, from the two outer bits, and the intermediate
physical-index sum acts on the inner bond. The Bell vectors together with
the two remaining coordinate vectors form the adapted basis used
throughout. Each weighted copy of a target contributes its weight raised
to the number of sites $N$, and weights are summed when they multiply the
same target; this is the power-sum mechanism of
Theorem~\ref{thm:asym_unblocked_coefficients}. The results below
establish tensor decompositions and periodic trace identities, not a
classification of length-dependent renormalization fixed points.

\begin{definition}[The idempotents of an adapted basis]
    \label{def:asymex_basis_proj}
    \lean{MPSTensor.basisProj, P6Compression.pairReorder, P6Compression.pairU,
        P6Compression.pairUinv, P6Compression.pairProjInt,
        P6Compression.pairBlockInt}
    \leanok
    \uses{prop:asym_projector_weighted_sum}
    For mutually inverse matrices $U$ and $U^{-1}$, define
    $\Pi_j=U^{-1}E_{jj}U$, the product of the $j$-th column of $U^{-1}$
    with the $j$-th row of $U$.
    Reorder a stacked bond quadruple into its two middle bits and its
    two outer bits. On the inner bond choose the adapted basis with
    first two vectors $(1,0,0,\pm1)$ and last two vectors the remaining
    coordinate vectors. The Bell idempotents carry the weighted targets.
    Their tensor products with target matrices are expressed in the
    original stacked coordinates by undoing the reordering.
\end{definition}

\begin{theorem}[The idempotents of an adapted basis are mutually orthogonal]
    \label{thm:asymex_basis_proj}
    \lean{MPSTensor.basisProj_mul_basisProj_self,
        MPSTensor.basisProj_mul_basisProj_of_ne, MPSTensor.mul_basisProj_mul}
    \leanok
    \uses{def:asymex_basis_proj}
    The idempotents satisfy $\Pi_j^2=\Pi_j$,
    $\Pi_j\Pi_k=0$ for $j\ne k$, and $U\Pi_jU^{-1}=E_{jj}$.
\end{theorem}
\begin{proof}
    \leanok
    Conjugation by $U$ returns $E_{jj}$, and products reduce to
    $E_{jj}E_{kk}=\delta_{jk}E_{jj}$.
\end{proof}

\begin{theorem}[From an integer matrix identity to a weighted direct sum]
    \label{thm:asymex_stacked_pair}
    \lean{P6Compression.stackedPair_letter_identity}
    \leanok
    \uses{def:asymex_basis_proj, prop:asym_projector_weighted_sum}
    Suppose $B=c_B\widehat B$ and $C_s=c_C\widehat C_s$, where the
    tensors with hats have integer entries and $s\in\{0,1\}$.
    Suppose a nonzero integer $n$ and integers $k_s$ satisfy, for each
    $s$ and each physical index $a$,
    \begin{align}
        n\,\mu_s\,\frac{c_C}{2} & =k_s\,c_B,
        \label{eq:asymex_stacked_pair_scalar}
    \end{align}
    \begin{align}
        n\,\widehat B^a & =\sum_sk_s\,(\widehat\Pi_s\otimes\widehat C_s^a),
        \label{eq:asymex_stacked_pair_letters}
    \end{align}
    where $\widehat\Pi_s=2\Pi_s$ and the right-hand side
    of~(\ref{eq:asymex_stacked_pair_letters}) is in stacked bond
    coordinates. Then $B^a=\sum_s\mu_s(\Pi_s\otimes C_s^a)$ in those
    coordinates, the hypothesis of
    Proposition~\ref{prop:asym_projector_weighted_sum}.
    Here $\mu_s$ denotes a scalar fusion weight, not a vertical
    multiplicity matrix.
\end{theorem}
\begin{proof}
    \leanok
    \uses{thm:asymex_basis_proj}
    Using the idempotents of Theorem~\ref{thm:asymex_basis_proj}, the
    $s$-th summand is
    $\mu_sc_C(\widehat\Pi_s\otimes\widehat C_s^a)/2$.
    Multiplication by $n$ and
    substitution of~(\ref{eq:asymex_stacked_pair_scalar})
    and~(\ref{eq:asymex_stacked_pair_letters}) give $nB^a$.
    Cancel the nonzero scalar $n$.
\end{proof}

\begin{theorem}[Normality from scaled matrix units]
    \label{thm:asymex_normal_units}
    \lean{P6Compression.isNormal_of_single_eq_smul}
    \leanok
    \uses{def:normal}
    Let a tensor be one nonzero scalar times an integer tensor. Suppose
    that for every bond position $(x,y)$ a letter of the integer tensor
    is a nonzero integer multiple of $E_{xy}$. Then its one-site
    matrices span the full bond matrix algebra, so it is normal at
    blocking length one in the sense of Definition~\ref{def:normal}.
    This establishes injectivity, not transfer spectral-radius
    normalization; a normal tensor (NT) in the RFP canonical form
    additionally requires that normalization.
\end{theorem}
\begin{proof}
    \leanok
    Dividing each designated letter by its nonzero factor exhibits every
    matrix unit in the one-site span.
\end{proof}

\section{A sign-weighted pair with a non-split extension}%
\label{sec:asymex_parity}

The tensors of this section are a construction of this development with
no literature source: a bond-dimension-two tensor $A$, the sign-weighted
pair of targets $(A,-A)$, and a bond-dimension-five source that carries
both targets together with a one-dimensional zero slot, in coordinates
where no direct-sum structure is visible. The construction tests three
features of Theorem~\ref{thm:asym_multiblock} on explicit data: the
remainder bound is attained
(Theorem~\ref{thm:asymex_parity_compression}), the periodic trace
identity acts as a length filter
(Theorem~\ref{thm:asymex_parity_trace}), and a word-level compression
need not come from sitewise intertwiners
(Theorem~\ref{thm:asymex_parity_no_intertwiner}).

Consider a two-state physical site with basis $\{\ket{0},\ket{1}\}$.
Define the bond-two MPS tensor by
\begin{align}
    A^0=P=\begin{pmatrix}1&0\\0&0\end{pmatrix},
    \qquad
    A^1=X=\begin{pmatrix}0&1\\1&0\end{pmatrix}.
    \label{eq:asymex_parity_letters}
\end{align}
For $N\geq1$, write its periodic amplitudes as
\begin{align}
    a_N(x)=\tr(A^{x_0}\cdots A^{x_{N-1}})
    \qquad (x\in\{0,1\}^N).
    \label{eq:asymex_parity_amplitudes}
\end{align}
The associated unnormalized state is
\begin{align}
    \ket{V^{(N)}(A)}
    =\sum_{x\in\{0,1\}^N}a_N(x)\ket{x},
    \label{eq:asymex_parity_state}
\end{align}
where $\ket{x}=\ket{x_0}\otimes\cdots\otimes\ket{x_{N-1}}$.
The state can be read off directly from its two letters: the matrix $P$
selects virtual coordinate $0$, while $X$ exchanges the two virtual
coordinates.

Multiplying both letters of $A$ by $-1$ gives one minus sign per
site, whatever the physical labels. For a word
$w=(i_1,\ldots,i_{|w|})$, write
$A^w=A^{i_1}\cdots A^{i_{|w|}}$, with the identity matrix as the
empty product. Then
\begin{align}
    (-A)^w=(-1)^{|w|}A^w.
    \label{eq:asymex_parity_weight}
\end{align}
The two weights $+1$ and $-1$ therefore produce a factor
$1+(-1)^N$ in the sum of their periodic amplitudes. The following
five-dimensional tensor realizes that sum without being a direct
sum in its displayed coordinates.

\begin{definition}[The graded pair and the five-dimensional source]
    \label{def:asymex_parity}
    \lean{ParityGraded.parA, ParityGraded.parTargets, ParityGraded.parB}
    \leanok
    \uses{def:mps_tensor}
    Let $A$ be the tensor in~(\ref{eq:asymex_parity_letters}).
    The two target tensors are $C_+=A$ and $C_-=-A$.
    The source is the bond-dimension-five tensor with letters
    \begin{align}
        B^0
         & =\begin{pmatrix}
                -1 & 1 & -1 & 1  & 3 \\
                2  & 2 & -2 & -1 & 0 \\
                1  & 2 & -2 & 0  & 2 \\
                1  & 3 & -3 & -1 & 2 \\
                -1 & 0 & 0  & 1  & 2
            \end{pmatrix}
        \label{eq:asymex_parity_source_zero}
    \end{align}
    and
    \begin{align}
        B^1
         & =\begin{pmatrix}
                3  & 1 & 0  & -2 & -3 \\
                -3 & 1 & -1 & 2  & 5  \\
                -1 & 2 & -1 & 1  & 3  \\
                0  & 1 & -1 & -1 & 1  \\
                2  & 1 & 0  & -1 & -2
            \end{pmatrix}.
        \label{eq:asymex_parity_source_one}
    \end{align}
\end{definition}

Neither source letter has a zero row or a zero column. Nevertheless,
a simultaneous change of bond coordinates reveals diagonal blocks
$A$, $-A$, and a one-dimensional zero block. The latter contributes
nothing to positive-length traces.

\begin{center}
    % Formula: the periodic trace tr(B^w)=(1+(-1)^{|w|})tr(A^w) of
    % Theorem~\ref{thm:asymex_parity_trace}; the source is not presented
    % as a literal stack of $A$ and $-A$, so the ring carries $B$ with
    % its coefficient, not a two-row stack.
    % Source: Definition~\ref{def:asymex_parity}.
    % Ink-to-index map: each node is the bond-dimension-five source
    % tensor $B$; the upper legs carry the letters of the word.
    % Contraction set: the horizontal virtual legs, closed around the
    % ring by the trace; no leg is left open.
    % Boundary signature: (virtual-west=0 | virtual-east=0 |
    % physical-up=3 | physical-down=0) for the finite schematic shown.
    % The physical legs remain open for the state; fixing their labels
    % gives the scalar coefficient tr(B^w). No virtual leg remains open.
    \tenkzkernel
    \begin{tenkz}[cols=4, boundary=periodic, physical=up]
        \tn[label pos=135, ports={90:physical:$i_1$}]{B} &
        \tn[label pos=135]{B} & \tn[skin=dots]{} &
        \tn[label pos=135, ports={90:physical:$i_N$}]{B}
    \end{tenkz}
    \qquad $\tr(B^w)$, computed in
    Theorem~\ref{thm:asymex_parity_trace}
\end{center}

\begin{theorem}[Normality of the graded block]
    \label{thm:asymex_parity_normal}
    \lean{ParityGraded.parA_isNormal}
    \leanok
    \uses{def:asymex_parity, def:normal}
    The block $A$ is normal: its length-two words span $\MN{2}$.
\end{theorem}

\begin{proof}
    \leanok
    Let $E_{ab}$ denote the matrix unit with its sole nonzero entry
    equal to one in row $a$, column $b$. Direct multiplication gives
    $E_{00}=A^0A^0$, $E_{01}=A^0A^1$, and $E_{10}=A^1A^0$.
    The remaining unit is $E_{11}=A^1A^1-A^0A^0$, so all four
    matrix units lie in the length-two word span.
\end{proof}

\subsection{Occupation parity and periodic amplitudes}

Call $n_1(x)=\sum_{v=0}^{N-1}x_v$ the occupation of the configuration $x$.
Let $z=\diag(1,-1)$ act on the physical site, and let
$Z=\diag(1,-1)$ act on the bond space. Although their matrices are
equal, these operators act on different spaces. The physical parity
operator has basis action
\begin{align}
    z^{\otimes N}\ket{x}=(-1)^{n_1(x)}\ket{x}.
    \label{eq:asymex_parity_basis}
\end{align}

\begin{lemma}[Occupation parity]
    \label{lem:asymex_parity_occupation}
    For each physical label $i\in\{0,1\}$, the tensor satisfies
    \begin{align}
        ZA^iZ=(-1)^iA^i.
        \label{eq:asymex_parity_covariance}
    \end{align}
    Consequently, $a_N(x)=0$ whenever $n_1(x)$ is odd, and
    \begin{align}
        z^{\otimes N}\ket{V^{(N)}(A)}=\ket{V^{(N)}(A)}.
        \label{eq:asymex_parity_invariance}
    \end{align}
\end{lemma}

\begin{proof}
    Direct multiplication gives $ZPZ=P$ and $ZXZ=-X$.
    Since $Z^2=\Id_2$, multiplying~(\ref{eq:asymex_parity_covariance})
    along the word yields
    \begin{align}
        Z(A^{x_0}\cdots A^{x_{N-1}})Z
        =(-1)^{n_1(x)}A^{x_0}\cdots A^{x_{N-1}}.
        \notag
    \end{align}
    Taking the trace and using cyclicity gives
    $a_N(x)=(-1)^{n_1(x)}a_N(x)$. Thus odd occupation forces
    $a_N(x)=0$. Applying~(\ref{eq:asymex_parity_basis}) to the
    expansion~(\ref{eq:asymex_parity_state}) proves
    (\ref{eq:asymex_parity_invariance}).
\end{proof}

The support of the state is further constrained by the relative positions of the
zeros.

\begin{lemma}[Amplitudes from cyclic runs]
    \label{lem:asymex_parity_runs}
    Suppose $x$ contains $m\geq1$ zeros. Read the sites cyclically, and
    let $r_1,\ldots,r_m$ be the numbers of ones between successive
    zeros. Define $b(r)=1$ for even $r$ and $b(r)=0$ for odd $r$.
    Then
    \begin{align}
        a_N(x)=\prod_{j=1}^{m}b(r_j).
        \label{eq:asymex_parity_runs}
    \end{align}
    For the word with no zeros,
    \begin{align}
        a_N(1,\ldots,1)=1+(-1)^N.
        \label{eq:asymex_parity_all_one}
    \end{align}
\end{lemma}

\begin{proof}
    The identities $X^2=\Id_2$ and $PXP=0$ imply
    $PX^rP=b(r)P$ and $\tr(PX^r)=b(r)$. Rotate the word to
    start at a zero, using cyclicity of the trace. Successive
    contractions give
    \begin{align}
        a_N(x)
         & =\tr(PX^{r_1}PX^{r_2}\cdots PX^{r_m}) \notag       \\
         & =\left(\prod_{j=1}^{m-1}b(r_j)\right)\tr(PX^{r_m})
        =\prod_{j=1}^{m}b(r_j).
        \notag
    \end{align}
    For the all-one word, $X^N=\Id_2$ when $N$ is even and
    $X^N=X$ when $N$ is odd. Their traces are $2$ and $0$,
    respectively, which proves~(\ref{eq:asymex_parity_all_one}).
\end{proof}

Thus a word containing a zero has amplitude one exactly when every
cyclic run of ones has even length. For example, $0101$ has even
occupation but zero amplitude, since both runs have length one.
The all-zero word has amplitude one at every positive length, so
$\ket{V^{(N)}(A)}$ is always nonzero in the specified computational basis.

\subsection{Asymmetric compression and length cancellation}

\begin{theorem}[Compression of the graded pair]
    \label{thm:asymex_parity_compression}
    % Principal declaration: ParityGraded.parityGraded_isReduction.
    \lean{ParityGraded.parityGraded_compression,
        ParityGraded.parityGraded_isReduction,
        ParityGraded.parityGraded_left_mul_right_of_ne,
        ParityGraded.parityGraded_dim_eq,
        ParityGraded.parityGraded_evalWord_remainder_eq_zero,
        ParityGraded.parityGraded_remainder_sq_ne_zero}
    \leanok
    \uses{def:asymex_parity, thm:asym_multiblock}
    The source of Definition~\ref{def:asymex_parity} carries the block
    triangular form of Theorem~\ref{thm:asym_multiblock} for
    $S=\{+,-\}$, with targets $C_+=A$, $C_-=-A$, one zero slot
    $z=1$, and dimension count $5=2+2+1$.
    There are compression maps $V_s\in\C^{5\times2}$ and
    $W_s\in\C^{2\times5}$ such that, for every $s,t\in S$,
    \begin{align}
        W_sV_t=\delta_{st}\Id_2.
        \label{eq:asymex_parity_biorthogonal}
    \end{align}
    For every word $w$ and each $s\in S$,
    \begin{align}
        W_sB^wV_s=C_s^w.
        \label{eq:asymex_parity_reduction}
    \end{align}
    Define the remainder letters by
    \begin{align}
        R^i=B^i-V_+A^iW_+-V_-(-A^i)W_-.
        \label{eq:asymex_parity_remainder}
    \end{align}
    Every product $R^w$ with $|w|\geq3$ vanishes, whereas
    $R^0R^0\neq0$. Thus the bound $|S|+z=3$ in clause (vi) of
    Theorem~\ref{thm:asym_multiblock} is attained.
\end{theorem}

\begin{proof}
    \leanok
    \uses{thm:asymex_explicit_gauge}
    Use the integer change of coordinates
    \begin{align}
        G=\begin{pmatrix}
              0  & -1 & 1  & 0  & 0  \\
              1  & 1  & 0  & -1 & -1 \\
              -1 & 0  & 0  & 1  & 1  \\
              1  & 0  & 0  & 0  & -1 \\
              0  & 1  & -1 & 0  & 1
          \end{pmatrix}.
        \label{eq:asymex_parity_gauge}
    \end{align}
    Its inverse is
    \begin{align}
        G^{-1}=\begin{pmatrix}
                   1 & 0 & 0 & 1 & 1 \\
                   0 & 1 & 1 & 0 & 0 \\
                   1 & 1 & 1 & 0 & 0 \\
                   0 & 0 & 1 & 1 & 0 \\
                   1 & 0 & 0 & 0 & 1
               \end{pmatrix}.
        \label{eq:asymex_parity_gauge_inv}
    \end{align}
    Direct multiplication verifies both inverse identities and gives
    \begin{align}
        GB^0G^{-1}
        =\begin{pmatrix}
             1 & 0 & 1  & 0 & 1 \\
             0 & 0 & 0  & 1 & 0 \\
             0 & 0 & -1 & 0 & 2 \\
             0 & 0 & 0  & 0 & 1 \\
             0 & 0 & 0  & 0 & 0
         \end{pmatrix}.
        \label{eq:asymex_parity_conj_zero}
    \end{align}
    For the second letter,
    \begin{align}
        GB^1G^{-1}
        =\begin{pmatrix}
             0 & 1 & 0  & 1  & 0 \\
             1 & 0 & 2  & 0  & 1 \\
             0 & 0 & 0  & -1 & 1 \\
             0 & 0 & -1 & 0  & 0 \\
             0 & 0 & 0  & 0  & 0
         \end{pmatrix}.
        \label{eq:asymex_parity_conj_one}
    \end{align}
    With block sizes $(2,2,1)$, the matrices in
    (\ref{eq:asymex_parity_conj_zero}) and
    (\ref{eq:asymex_parity_conj_one}) are block upper triangular,
    with diagonal blocks $A^i$, $-A^i$, and $0$.
    As in Theorem~\ref{thm:asymex_explicit_gauge}, take $W_+$ and
    $W_-$ to be rows $(0,1)$ and $(2,3)$ of $G$, respectively,
    and $V_+$ and $V_-$ to be the corresponding columns of $G^{-1}$.
    The inverse identities give~(\ref{eq:asymex_parity_biorthogonal}).
    Diagonal blocks multiply under word evaluation, giving
    (\ref{eq:asymex_parity_reduction}).

    Subtracting the two diagonal target blocks leaves a strictly
    upper triangular three-block remainder in these coordinates.
    The compression remainder bound therefore gives $R^w=0$
    for $|w|\geq3$. To check sharpness in the original coordinates,
    substitution into~(\ref{eq:asymex_parity_remainder}) gives
    \begin{align}
        R^0=\begin{pmatrix}
                -1 & 2 & -2 & 1 & 3 \\
                1  & 2 & -2 & 0 & 1 \\
                0  & 3 & -3 & 1 & 3 \\
                0  & 3 & -3 & 0 & 3 \\
                -1 & 1 & -1 & 1 & 2
            \end{pmatrix}.
        \notag
    \end{align}
    Its square has $(0,1)$ entry $2$, so $R^0R^0\neq0$.
\end{proof}

\begin{theorem}[Periodic trace of the sign-weighted pair]
    \label{thm:asymex_parity_trace}
    \lean{ParityGraded.parityGraded_trace_evalWord}
    \leanok
    \uses{def:asymex_parity}
    For every nonempty word $w$,
    \begin{align}
        \tr(B^w)=(1+(-1)^{|w|})\tr(A^w).
        \label{eq:asymex_parity_trace}
    \end{align}
\end{theorem}

\begin{proof}
    \leanok
    \uses{thm:asymex_parity_compression, prop:asym_compression_trace}
    Proposition~\ref{prop:asym_compression_trace} applied to the datum of
    Theorem~\ref{thm:asymex_parity_compression} gives
    $\tr(B^w)=\tr(A^w)+\tr((-A)^w)$. By
    (\ref{eq:asymex_parity_weight}), the second summand is
    $(-1)^{|w|}\tr(A^w)$, proving~(\ref{eq:asymex_parity_trace}).
\end{proof}

Define $\ket{V^{(N)}(B)}$ by the same periodic contraction as
in~(\ref{eq:asymex_parity_state}), with $A$ replaced by $B$.
Applying~(\ref{eq:asymex_parity_trace}) to each basis coefficient
gives, for $N\geq1$,
\begin{align}
    \ket{V^{(N)}(B)}
    =(1+(-1)^N)\ket{V^{(N)}(A)}.
    \label{eq:asymex_parity_length}
\end{align}
The source acts as a length filter: at odd length the periodic vector
vanishes, while at even length every amplitude of $A$ is doubled.
For instance, at odd length the all-zero amplitude of $A$ is one but its
amplitude in $B$ vanishes; at even length the run restrictions of
Lemma~\ref{lem:asymex_parity_runs} apply with doubled amplitudes.

The positive-length hypothesis is necessary. The empty-word
traces are $\tr(\Id_5)=5$ and $\tr(\Id_2)=2$, whereas substituting
$|w|=0$ into~(\ref{eq:asymex_parity_trace}) would give $5=4$.
The missing contribution is the empty trace of the
one-dimensional zero block, which is one; its traces vanish only
at positive length.

\begin{theorem}[The weight $-1$ block has no sitewise intertwiner]
    \label{thm:asymex_parity_no_intertwiner}
    \lean{ParityGraded.parNeg_right_intertwiner_eq_zero,
        ParityGraded.parNeg_left_intertwiner_eq_zero}
    \leanok
    \uses{def:asymex_parity}
    If $V\in\C^{5\times2}$ satisfies $B^iV=V(-A^i)$ for both letters, then
    $V=0$; if $W\in\C^{2\times5}$ satisfies $WB^i=(-A^i)W$ for both
    letters, then $W=0$.
\end{theorem}

\begin{proof}
    \leanok
    Each condition is a homogeneous linear system in ten unknown
    entries. For the right intertwiner, select entries
    $(3,0)$, $(4,0)$, $(2,1)$, $(3,1)$, and $(4,1)$ of
    $B^0V=V(-A^0)$, and all five entries of column $0$ of
    $B^1V=V(-A^1)$. Rational linear combinations of these ten
    equations force all entries of column $0$ of $V$ to vanish.
    The selected equations for letter $1$ then force column $1$
    to vanish.

    For the left intertwiner, select entries $(0,4)$, $(1,0)$,
    $(1,2)$, $(1,3)$, and $(1,4)$ of $WB^0=(-A^0)W$, and all
    five entries of row $0$ of $WB^1=(-A^1)W$. Rational linear
    combinations force row $0$ of $W$ to vanish, after which the
    letter-$1$ equations force row $1$ to vanish.
\end{proof}

The nonzero maps $W_-$ and $V_-$ in
(\ref{eq:asymex_parity_reduction}) therefore give a word-level
compression onto $-A$ without either map being a sitewise
intertwiner. The trace identity and the sharp remainder bound
hold despite this obstruction.
